\documentclass[12pt]{article}
\usepackage[T1]{fontenc}
\usepackage[utf8]{inputenc}
\usepackage{lmodern}
\usepackage{textcomp}
\usepackage{enumitem}
\usepackage{geometry}
\usepackage{graphicx}
\usepackage{amsmath, amssymb}
\geometry{margin=1in}
\begin{document}
\begin{center}
Sociology - Graduate Level Assessment 22 \
Total Marks: 40 \
Answer all questions.
\end{center}

\section*{Multiple Choice (10 marks)}
\begin{enumerate}[label=\arabic*.]
\item Disciplinary power was exercised in the 19 th century factories by:
\begin{enumerate}[label=(\alph*)]
    \item the use of corporal punishment by employers
    \item excluding women from participating in waged labour
    \item punctuality, uninterrupted work and the threat of dismissal
    \item making routine tasks less monotonous
\end{enumerate}
\item Environmentalist social movements are global in the sense that:
\begin{enumerate}[label=(\alph*)]
    \item they increase our awareness of risks that affect the whole planet
    \item they appeal to universal values and human rights
    \item they use global media to generate publicity
    \item all of the above
\end{enumerate}
\item Which of the following did Domhoff not identify as a process of decision-making in the USA?
\begin{enumerate}[label=(\alph*)]
    \item the ideology process
    \item the exploitation process
    \item the policy-formation process
    \item the candidate-selection process
\end{enumerate}
\item Which of the following is not identified by Fulcher \& Scott as a criterion of community?
\begin{enumerate}[label=(\alph*)]
    \item a shared sense of identity and belonging together
    \item common activities involving all-round relationships
    \item a fixed geographical location
    \item collective action based on common interests
\end{enumerate}
\item Marx (1844) claimed that capitalism had alienated workers from:
\begin{enumerate}[label=(\alph*)]
    \item the product they were making
    \item the production process
    \item each other and humanity in general
    \item all of the above
\end{enumerate}
\end{enumerate}

\section*{True / False (10 marks)}
\textit{Answer True or False}
\begin{enumerate}[label=\arabic*.]
\setcounter{enumi}{5}
\item True or False: 'Snowballing' is a technique used in probability sampling to ensure a representative sample across diverse population segments.
\item True or False: The shift to joint parenthood after divorce requires both parents to cooperate in the ongoing care of their children.
\item True or False: Seven people waiting for a bus at a bus stop exemplify an aggregate, as they share a temporary physical space without social interaction.
\item True or False: Weber categorized effective social action as one of the four ideal types of social action in his analysis.
\item True or False: Technological surveillance has primarily increased personal privacy for consumers and employees in the workplace.
\end{enumerate}

\section*{Long Form Response (20 marks)}
\textit{Answer each question with a clear and complete explanation.}
\begin{enumerate}[label=\arabic*.]
\setcounter{enumi}{10}
\item Discuss the significance of personal documents, such as a household account book, in sociological research, particularly focusing on their closed access and implications for data analysis.
\item Discuss Blauner's (1964) analysis of alienation in the workplace, specifically focusing on car manufacturing in assembly plants. What factors contribute to this alienation, and what are its implications for employees?
\end{enumerate}

\vfill
\noindent\textit{End of Paper}
\end{document}